\documentclass[12pt]{article}
\usepackage[T1]{fontenc}
\usepackage[utf8]{inputenc}
\usepackage{lmodern}
\usepackage{textcomp}
\usepackage{enumitem}
\usepackage{geometry}
\geometry{margin=1in}
\begin{document}
\begin{center}
Answer Key - Biology - Graduate Level Assessment 4 \
\small (Answers are ordered exactly as in the question paper.)
\end{center}

\section*{Multiple Choice}
\begin{enumerate}[label=\arabic*.]
\item \textbf{Answer:} C
\\ \textit{Options:} A) Bundle sheath, companion cell, phloem parenchyma, fibers; B) Phloem parenchyma, bundle sheath, fibers, companion cell; C) Bundle sheath, phloem parenchyma, companion cell, sieve tube; D) Bundle sheath, fibers, companion cell, sieve tube; E) Fibers, bundle sheath, companion cell, sieve tube
\\ \textit{Note:} The correct answer outlines the symplastic pathway for sucrose movement, as it describes the route through which sucrose is transferred from mesophyll cells through plasmodesmata into the bundle sheath, then to phloem parenchyma, followed by companion cells, and finally into the sieve tube, illustrating the interconnected plant cell structures facilitating this transport.
\item \textbf{Answer:} A
\\ \textit{Options:} A) The thick muscular wall of the stomach; B) Antibacterial properties of stomach acid; C) Gastric juice; D) High pH levels in the stomach; E) Presence of beneficial bacteria
\\ \textit{Note:} The thick muscular wall of the stomach provides structural integrity and protection, preventing the digestive acids and enzymes from damaging the stomach tissue itself.
\item \textbf{Answer:} A
\\ \textit{Options:} A) Shivering is a method the body uses to relieve pain.; B) Shivering is a tactic to ward off predators by appearing agitated.; C) Shivering is a way to burn excess calories; D) Shivering is a response to fear or anxiety; E) Shivering is an involuntary action to prepare muscles for physical activity.
\\ \textit{Note:} Shivering is not primarily a method to relieve pain; rather, it is a physiological response to cold that generates heat through involuntary muscle contractions to maintain core body temperature.
\item \textbf{Answer:} A
\\ \textit{Options:} A) Taste, wind direction, humidity; B) Atmospheric oxygen levels, plant growth, tidal patterns; C) Infrared radiation, sound frequency, barometric pressure; D) Moon phases, echo location, gravitational pull; E) Sunlight intensity, water current, soil composition
\\ \textit{Note:} Animals utilize various environmental cues, such as taste, wind direction, and humidity, to navigate during migration, as these factors provide critical information about their surroundings and help them orient themselves effectively.
\item \textbf{Answer:} A
\\ \textit{Options:} A) Without HCl, more bacteria would be killed in the stomach.; B) No HCl would lead to reduced mucus production in the stomach lining.; C) The absence of HCl would cause the stomach to digest its own tissues.; D) No HCl in the stomach would significantly increase the pH of the small intestine.; E) The absence of HCl would speed up the digestive process.
\\ \textit{Note:} The correct answer is based on the fact that hydrochloric acid (HCl) in the stomach plays a crucial role in creating an acidic environment that helps to kill harmful bacteria and pathogens ingested with food, thus preventing infections and supporting overall digestive health.
\end{enumerate}

\section*{True / False}
\begin{enumerate}[label=\arabic*.]
\setcounter{enumi}{5}
\item \textbf{Answer:} True
\\ \textit{Options:} A) True; B) False
\\ \textit{Note:} Greater anterior laxity of the uninjured knee was associated with poorer stability and functional outcomes after ACL reconstruction. Excessive anterior laxity of the uninjured knee thus appears to represent a risk factor for inferior outcomes.
\item \textbf{Answer:} False
\\ \textit{Options:} A) True; B) False
\\ \textit{Note:} The survey shows significant dissatisfaction amongst consultant radiologists with the current service, confirms a low number of paediatric radiologists taking on this work, and suggests the potential to increase numbers of radiology child abuse experts by 27\% if given improved training and support. Appropriate service and education strategies should be implemented.
\item \textbf{Answer:} False
\\ \textit{Options:} A) True; B) False
\\ \textit{Note:} Early oral intake is possible after laparotomy and colorectal resection. Thus, the laparoscopic surgeon's claim of early tolerated oral intake may not be unique to laparoscopy.
\item \textbf{Answer:} True
\\ \textit{Options:} A) True; B) False
\\ \textit{Note:} A higher PV in PMR increases the risk of prolonged steroid therapy and late GCA. Female sex and particular HLA alleles may increase the risk of late GCA. Starting patients on\textgreater{}15 mg prednisolone is associated with a prolonged steroid duration.
\item \textbf{Answer:} False
\\ \textit{Options:} A) True; B) False
\\ \textit{Note:} The incidence of venous thrombosis during hospitalisation in a department of general internal medicine is low and does not justify prophylaxis in all internal patients. Cancer is a strong risk factor for hospital-acquired thrombosis in the medical ward. Further studies may answer the question as to whether thrombosis prophylaxis in this subgroup is feasible.
\end{enumerate}

\section*{Long Form Response}
\begin{enumerate}[label=\arabic*.]
\setcounter{enumi}{10}
\item \textbf{Answer:} The use of human subjects in inheritance studies presents significant challenges due to the intricate interplay of genetic complexity and environmental factors. Human genetics is characterized by complex gene interactions, where multiple genes can influence a single trait, complicating the attribution of inheritance patterns. Additionally, environmental factors can modify phenotypic expressions, making it difficult to isolate genetic contributions. Ethical considerations further complicate these studies, as researchers must navigate issues related to consent, privacy, and potential harm to participants. Furthermore, demographic limitations, such as small family sizes and the extended lifespan of humans, restrict the sample sizes available for study, reducing the statistical power and generalizability of findings. Together, these factors create a multifaceted landscape that challenges the design and interpretation of inheritance studies in human populations.
\\ \textit{Note:} The answer accurately addresses the multifaceted challenges of using human subjects in inheritance studies by emphasizing the complexities of gene interactions, the influence of environmental factors on phenotypic expressions, and the ethical dilemmas that arise in such research.
\item \textbf{Answer:} Linkage disequilibrium (LD) in populations is promoted by several factors, including genetic drift, selection, population structure, and the presence of physical linkage between alleles on the same chromosome. Genetic drift can cause certain allele combinations to become more prevalent over generations, while selection may favor specific haplotypes, further enhancing LD. Population structure, such as subpopulation differentiation and limited gene flow, can also lead to non-random associations between alleles. Contrary to these mechanisms, random mating does not contribute to linkage disequilibrium because it promotes the mixing of alleles across generations, allowing for the independent assortment of alleles and thereby breaking down any existing associations. Thus, random mating serves to maintain or restore equilibrium rather than foster disequilibrium.
\\ \textit{Note:} The answer correctly identifies that factors such as genetic drift, selection, and population structure contribute to linkage disequilibrium by creating non-random associations between alleles, while also noting that random mating does not contribute to LD because it promotes allele mixing and thus breaks down these associations.
\end{enumerate}

\vfill
\noindent\textit{End of Answer Key}
\end{document}